(1) 入射场可写 \(p_i=A_i\cos(k_ix-\omega_it)\)，其中 \(k_i=\omega_i/c\)。反射场写 \(p_r=A_r\cos(-k_rx-\omega_rt+\phi)\)，其中正数 \(k_r=\omega_r/c\)。活塞轨迹 \(x=Ut\) 上的两个振荡相位必须有相同时间变化率。

(2) 壁面相位变化率，入射为 \(k_iU-\omega_i\)，反射为 \(-k_rU-\omega_r\)。两者相等给 \(\omega_i(1-U/c)=\omega_r(1+U/c)\)，故 \(\boxed{\omega_r=\omega_i(c-U)/(c+U)}\)，反射波长 \(\boxed{\lambda_r=2\pi c/\omega_r=\lambda_i(c+U)/(c-U)}\)。这是两次移动边界Doppler效应的非相对论声学形式。

(3) 若 \(U<0\)（壁向波源靠近），频率比 \((c-U)/(c+U)>1\)，出现蓝移，前提仍为 \(|U|<c\)。若 \(U=0\)，则 \(\omega_r=\omega_i\)。当 \(U\to c^-\)，\(\omega_r\to0\)，波长趋无穷；运动边界已近于入射波前速度，线性局域稳态近似会特别敏感。

(4) 移动壁会对流体做正功或从流体提取功，反射过程中动量方向反转，故活塞受到非零时间平均辐射压强。固定壁无需外界机械功，入射与反射能流相等；运动壁则有 \(P_{\rm acoustic,out}-P_{\rm acoustic,in}=P_{\rm wall}\) 的机械功账本，应结合移动边界振幅关系决定具体大小，不能仅凭频率比硬套振幅恒定。
